From~\eqref{equ:kp}, we know that labeled posets $(P,\om)$ and $(Q,\om')$  satisfy 
\[
\kpw K_{(Q,\om')} = K_{(P,\om) \sqcup (Q,\om')}.
\]
Thus we can interpret the product $F_\alpha F_\beta$ as a $(P,\om)$-partition enumerator for a disjoint union of two appropriately labeled chains.  One fact that we need later from this interpretation is summarized in the following lemma.

\begin{lemma}\label{lem:deg2}
Let $f_1$ and $f_2$ be two non-constant elements of \qsym.
Then the $F$-support of $f_1 f_2$ contains at least one composition $\gamma$ whose first part $\gamma_1$ is at least $2$.
The same is true when focusing on the last part of the compositions appearing in the $F$-support of $f_1 f_2$.
\end{lemma}
\begin{proof}
First consider two non-empty compositions $\alpha$ and $\beta$.  Construct two labeled chains $(P,\om)$ and $(Q,\om')$ such that $F_\alpha=K_{(P,\om)}$ and $F_\beta=K_{(Q,\om')}$.  Thus $F_\alpha F_\beta = K_{(P,\om) \sqcup (Q,\om')}$.  It is obvious that there is at least one element $\pi$ in the set $\L((P,\om) \sqcup (Q,\om'))$ with $\pi_1<\pi_2$.  By Theorem~\ref{theo:KinF}, this gives rise to an element $F_\gamma$ in the $F$-support of $F_\alpha F_\beta$ with $\gamma_1\ge2$.

Now, let us consider two non-constant elements $f_1$ and $f_2$ of \qsym.
Let us denote by $\alpha$ (\resp $\beta$) the lexicographically maximal composition in the $F$-support of $f_1$ (\resp $f_2$).
It is clear that the lexicographically maximal element in the product $f_1 f_2$ comes from the product $F_\alpha F_\beta$,
thus the case above applies.

The case of the last part is similar, requiring just two tweaks.  First, if $\ell$ denotes the number of elements in $(P,\om) \sqcup (Q,\om')$, it is obvious that there is at least one element $\pi$ in the set $\L((P,\om) \sqcup (Q,\om'))$ with $\pi_{\ell-1}<\pi_\ell$.   This yields an element $F_\gamma$ in the $F$-support of $F_\alpha F_\beta$ such that the last part of $\gamma$ is at least 2. Secondly, instead of considering the lexicographically maximal composition, we use the composition whose reversal is lexicographically maximal.
\end{proof}
